\lysh{15823}{4}
\begin{alist}
\item Από την ταυτότητα της διαίρεσης έχουμε $P(x) = (4x^2 - 1)(3x - 2) + 1$, οπότε η ζητούμενη εξίσωση γίνεται ισοδύναμα
\[P(x) = 1 \Leftrightarrow (4x^2 - 1)(3x - 2) + 1 = 1 \Leftrightarrow (4x^2 - 1)(3x - 2) = 0\]
Εύκολα πλέον βρίσκουμε ότι οι λύσεις της εξίσωσης είναι οι αριθμοί $\pm\dfrac{1}{2}$ και $\dfrac{2}{3}$.
\item Δείξαμε στο α) ότι αριθμοί που ικανοποιούν την εξίσωση $P(x) = 1$ είναι οι $\pm\dfrac{1}{2}$ και $\dfrac{2}{3}$. \\
Συνεπώς για να ήταν $P(\log 5) = 1$, θα έπρεπε $\log 5 = \dfrac{1}{2}$ ή $\log 5 = -\dfrac{1}{2}$ ή $\log 5 = \dfrac{2}{3}$. \\
Όμως $\log 5 \neq -\dfrac{1}{2}$ αφού $5 > 1 \Leftrightarrow \log 5 > 0$. \\
Επίσης $\log 5 = \dfrac{1}{2} \Leftrightarrow 5 = 10^{\frac{1}{2}} \Leftrightarrow 5 = \sqrt{10}$ το οποίο είναι άτοπο. \\
Τέλος $\log 5 = \dfrac{2}{3} \Leftrightarrow 5 = 10^{\frac{2}{3}} \Leftrightarrow 5 = \sqrt[3]{10^2}$ το οποίο επίσης είναι άτοπο. \\
Συνεπώς $P(\log 5) \neq 1$.
\item Είναι $P(x) = (4x^2 - 1)(3x - 2) + 1$, οπότε $P(-1) = (4(-1)^2 - 1)(3(-1) - 2) + 1 = -14$ και $P(0) = (-1)(-2) + 1 = 3$. \\
Αφού οι τιμές $P(-1)$, $P(0)$ είναι ετερόσημες, η εξίσωση $P(x) = 0$ έχει μία τουλάχιστον ρίζα στο $(-1, 0)$.
\end{alist}

